\begin{center}
	\Large{\textbf{EXECUTIVE SUMMARY}}
\end{center}
\vspace{1em}
\normalsize

\hspace{1cm} Cybersecurity threats have grown increasingly sophisticated, with automated bots, script-based exploit tools, and advanced persistent threat actors continuously probing networked systems for vulnerabilities. Conventional honeypot solutions are limited in their ability to respond intelligently to such attacks, as they record network interactions passively without adapting their behaviour or extracting structured threat intelligence from attacker sessions. This final-year engineering project presents the design, implementation, and validation of a \textbf{Behaviour-Aware Honeypot} — an active deception system that integrates machine learning, natural language generation, and stateful behavioural analysis to engage a wide range of adversarial profiles. The system deploys two externally accessible trap services: a Secure Shell (SSH) honeypot that accepts any credential and presents a convincing fake shell environment, and an HTTP honeypot styled as a fictional corporate web application with deliberately seeded vulnerabilities, both co-ordinated by a central REST API and isolated within Docker containers. Attacker interactions are classified by a cascaded machine-learning pipeline comprising a TF-IDF Random Forest fast-path classifier, a DistilBART zero-shot deep classifier, and a PyTorch bidirectional Long Short-Term Memory (BiLSTM) sequence profiler. A deterministic state machine accumulates a per-attacker risk score across both protocols, enabling cross-service kill-chain detection. Upon confirming a threat, the system generates contextually consistent terminal responses using the Phi-3-mini large language model executed locally via Ollama, ensuring no attacker data is transmitted to external services. Forensic integrity is guaranteed through a SHA-256 hash-chained event log, and a live Streamlit dashboard provides the operator with real-time visibility and AI-mode control. Systematic testing across four adversarial profiles and an independent red-team evaluation confirmed that all performance targets were met: deterministic shell commands completed in under 100\,ms, attack classification achieved maximum confidence for canonical payloads, and twelve MITRE ATT\&CK techniques were correctly mapped across sessions. The project demonstrates that AI-driven adaptive deception is a viable and deployable approach to modern honeypot engineering.
